\bookchapter{Capítulo 10: Pagos móviles y servicios cotidianos}

\chapterintro{Sinceramente, la generación de mis padres al principio se resistía a “pagar con el teléfono”. “Si el dinero está en el teléfono, ¿qué pasa si te lo roban?”. Es un temor razonable, pero la realidad es que la probabilidad de perder el efectivo del bolsillo es mucho mayor que la de que te clonen un pago móvil. El pago móvil no consiste en guardar el dinero en el teléfono; es solo una llave, una llave que abre tu cuenta bancaria.}

\booksection{I. ¿Cómo se vincula una tarjeta bancaria?}

Tanto si usas Alipay como WeChat, la primera vez que vincules una tarjeta debes seguir estos pasos:

\begin{enumerate}
  \item Abre la App, busca "Mi" o "Cartera". Pulsa "Añadir tarjeta bancaria".
  \item Introduce el número de tarjeta. La mayoría de las App permiten ahora escanear la tarjeta con la cámara, sin teclear dígito a dígito.
  \item El banco enviará un SMS con un código de verificación al número que tengas registrado. Introdúcelo y define una contraseña de pago de 6 dígitos.
  \item Esta contraseña de pago debería ser distinta de la contraseña de la tarjeta bancaria y de la de bloqueo de pantalla; no las iguales por comodidad.
\end{enumerate}

Un truco: si te preocupa la seguridad, puedes abrir en la App del banco una cuenta electrónica de categoría II o III dedicada a vincular al pago móvil. Mete en esa tarjeta solo el dinero de gastos diarios; si pasa algo, la pérdida es limitada.

\booksection{II. El código de pago es “quien lo ve, cobra”}

Mucha gente confunde el código de pago con el de cobro; la diferencia es importante.

\begin{itemize}
  \item \textbf{Código de pago}: el código de barras + código QR que muestras al cajero para que lo escanee. Equivale a tu "cartera electrónica": quien lo escanea puede retirar dinero de tu cuenta. Por eso\textbf{nunca lo captures ni lo envíes a nadie}. Este código se renueva automáticamente cada minuto, precisamente para evitar usos fraudulentos.
  \item \textbf{Código de cobro}: otros lo escanean para transferirte dinero. Sí se puede capturar y guardar, y se puede enviar a otros; es seguro.
\end{itemize}

Ten clara esta diferencia y evitarás muchas trampas de fraude. Cualquiera que te pida “envíame una captura del código de pago” es, al 100 \%, un estafador.

\booksection{III. Estos ajustes de seguridad, 5 minutos que ahorran muchos dolores de cabeza}

\begin{itemize}
  \item \textbf{Pago con huella/rostro}: en los ajustes de pago de Alipay y WeChat, activa el pago biométrico. Así, cada pago exige huella o rostro; aunque pierdas el teléfono, quien lo encuentre no podrá gastar tu dinero.
  \item \textbf{Bloqueo de seguridad}: Alipay y WeChat tienen una función de "bloqueo de seguridad". Al activarla, entrar en la página "Cartera" de la App exige una verificación adicional. Ruta: Alipay Mi Ajustes Cuenta y seguridad Centro de seguridad Bloqueo de seguridad.
  \item \textbf{Revisión del pago sin contraseña}: mucha gente ha suscrito varios pagos sin contraseña y hasta se ha olvidado. En los "ajustes de pago" de Alipay y WeChat, busca "Pago sin contraseña/cobro automático" y revisa qué servicios tienes suscritos; cancela los que no necesites.
  \item \textbf{Gestión de dispositivos conectados}: revisa periódicamente qué dispositivos han iniciado sesión en tu cuenta. Si ves uno desconocido, cierra la sesión de inmediato. Ruta: Alipay Ajustes Cuenta y seguridad Gestión de dispositivos.
\end{itemize}

\booksection{IV. Autobús y metro: el teléfono es la tarjeta de transporte}

Hoy en día, son pocos los que usan tarjeta de transporte física en el día a día. Para pasar el teléfono por el torniquete hay dos caminos:

\begin{itemize}
  \item \textbf{Tarjeta de transporte virtual en la cartera del teléfono}: abre una tarjeta de transporte virtual directamente en Apple Wallet, Huawei Wallet o Xiaomi Wallet. La ventaja: no hace falta desbloquear el teléfono para pasar; basta con acercarlo, la respuesta es la más rápida y funciona sin red. En algunas ciudades, la tarjeta virtual ofrece las mismás ventajas de transbordo que la física.
  \item \textbf{Código de viaje de Alipay/WeChat}: busca "código de viaje" en la App y activa el código de autobús o metro de tu ciudad. La ventaja: no hay que recargar; se descuenta del saldo o de la tarjeta bancaria vinculada. Pero requiere conexión y, antes de pasar, hay que abrir la App y encontrar el código.
\end{itemize}

Personalmente recomiendo más la tarjeta de transporte virtual de la cartera del teléfono. Aunque la primera vez haya que cargar 10 o 20 yuanes, es mucho más cómoda, sobre todo cuando se corre al metro y no hay que pararse a desbloquear el teléfono y abrir la App.

\booksection{V. Pagos de suministros, comida a domicilio, taxi}

\begin{itemize}
  \item \textbf{Pago de agua, electricidad y gas}: en "Pagos de suministros" de Alipay o WeChat, vincula los números de cuenta de tu casa de agua, electricidad y gas. Cada mes, al recibir la factura, entra y paga. El número de cuenta suele estar en la factura en papel o se puede consultar por teléfono al servicio de atención al cliente.
  \item \textbf{Comparar precios de comida a domicilio}: el mismo restaurante suele tener precios y ofertas distintas en Meituan y en Ele.me. Antes de pedir, mira en ambas; a veces te ahorras unos yuanes. Además, la función "Pin Hao Fan" de Meituan, pidiendo con amigos, sale bastante más barato.
  \item \textbf{Seguridad del taxi}: en DiDi o en el taxi de Gaode, recuerda configurar "contacto de emergencia" y "compartir trayecto". Al pedir un taxi, la App comparte automáticamente tu trayecto con el contacto que hayas definido. Además, de noche prioriza el "taxi" frente al "coche exprés"; el taxi tiene regulación del sector y mayor seguridad.
\end{itemize}

\booksection{VI. Gestíona con el teléfono la seguridad social, la sanidad y los trámites}

El cambio en los últimos años ha sido grande; muchos trámites que antes requerían ir a una oficina y hacer cola se resuelven ahora en unos minutos desde el teléfono.

\begin{itemize}
  \item \textbf{Certificado electrónico de seguro médico}: busca "certificado electrónico de seguro médico" en Alipay o WeChat; tras activarlo, puedes acudir al médico o comprar medicamentos en la farmacia pasando el teléfono, sin llevar la tarjeta física. La primera vez exige verificación facial; después, cada uso genera un código dinámico.
  \item \textbf{Consulta de la seguridad social}: descarga la App "12333" o usa el "Centro ciudadano" de Alipay para consultar el historial de cotizaciones, el saldo de la cuenta de pensiones o solicitar el subsidio por desempleo.
  \item \textbf{IRPF}: la liquidación anual del impuesto sobre la renta, de marzo a junio, se puede hacer en la App "Impuesto sobre la renta personal". Recuerda revisar si tienes deducciones especiales aplicables (educación de los hijos, intereses de hipoteca, cuidado de mayores, etc.); todas te ayudan a pagar menos o a recibir devolución.
\end{itemize}

\begin{tcbwarn}{Recuerda esto para evitar fraudes}
\begin{itemize}
  \item La policía, los bancos y el servicio al cliente nunca te pedirán que transfieras a una "cuenta segura".
  \item Cualquiera que te pida enviarle una captura del código de pago es un estafador.
  \item No hagas clic en enlaces desconocidos, sobre todo los de SMS del tipo "ETC caducado" o "tarjeta de la seguridad social suspendida".
  \item Los reembolsos de compras online se devuelven por la misma vía; no hace falta que fácilites el número de tu tarjeta bancaria.
\end{itemize}

\end{tcbwarn}

\begin{tcbpractice}{Hazlo ahora: revisa la seguridad de tus App de pago}
\begin{enumerate}
  \item Abre Alipay y WeChat y confirma que tienes activado el pago con huella/rostro.
  \item Encuentra la función "Bloqueo de seguridad" y actívala.
  \item Revisa la lista de "Pago sin contraseña/cobro automático" y desactiva los servicios que ya no necesites.
  \item Revisa la lista de "Dispositivos conectados" y elimina los que no reconozcas.
  \item Activa tu certificado electrónico de seguro médico; la próxima vez que vayas al médico no hará falta llevar la tarjeta física.
\end{enumerate}

\end{tcbpractice}
